\setcounter{section}{28}

  \zwischenueberschrift{Krümmung}

Es seien
\mathkor {} {V} {und} {W} {}
zweifach stetig differenzierbare Vektorfelder auf einer differenzierbaren Mannigfaltigkeit $M$ und sei $\nabla$ ein
\definitionsverweis {linearer Zusammenhang}{}{}
auf einem
\definitionsverweis {Vektorbündel}{}{}
\maabb {p} { E } { M
} {}
über $M$. Es ist
\maabbeledisp {\nabla_W} {C^2(E) } { C^1 (E)
} { s } { \nabla_W (s)
} {,}
die
\definitionsverweis {vertikale Ableitung}{}{.}
Auf das Ergebnis $\nabla_W (s)$ kann man wiederum die vertikale Ableitung in Richtung $V$ anwenden und erhält den stetigen Schnitt
\mathl{\nabla_V  \left(  \nabla_W (s)  \right)}{.} Wenn alle Objekte $C^\infty$ sind, muss man sich über die Differenzierbarkeitsgrade keine Gedanken machen. Hier betrachten wir den Ausdruck

\mavergleichskettedisp
{\vergleichskette
{ R(V,W)
}
{ =} { \nabla_V \nabla_W - \nabla_W \nabla_V- \nabla_{[V,W]}
}
{ } {
}
{ } {
}
{ } {
}
}
{}{}{,}
der aus
\zusatzklammer {hinreichend differenzierbaren} {} {}
Schnitten in $E$ wieder einen Schnitt in $E$ produziert. Dabei ist
\mathl{[V,W]}{} die
\definitionsverweis {Lie-Klammer}{}{}
von Vektorfeldern.

\inputdefinition
{}
{

Es sei $M$ eine zweifach stetige
\definitionsverweis {differenzierbare Mannigfaltigkeit}{}{}
und sei
\maabb {} { E } { M
} {}
ein zweifach stetig
\definitionsverweis {differenzierbares Vektorbündel}{}{}
über $M$ mit einem
\definitionsverweis {linearen Zusammenhang}{}{}
$\nabla$. Dann nennt man zu differenzierbaren
\definitionsverweis {Vektorfeldern}{}{}
\mathkor {} {V} {und} {W} {}
auf $M$ die Abbildung
\maabbeledisp {} {C^2(E)} {C^0(E)
} { s } { \nabla_V \nabla_W (s) - \nabla_W \nabla_V (s) - \nabla_{[V,W]}(s)
} {,}
den
\definitionswort {Krümmungsoperator}{}
zu $V$ und $W$. Er wird mit $R(V,W)$ bezeichnet.

}

__NOEDITSECTION__

\inputfaktbeweis
{Reelles Vektorbündel/Trivial/Linearer Zusammenhang/Krümmungsoperator/Beschreibung/Fakt}
{Lemma}
{}
{

\faktsituation {Es sei

\mavergleichskette
{\vergleichskette
{ U
}
{ \subseteq }{ \R^d
}
{  }{
}
{    }{
}
{   }{
}
}
{}{}{}
\definitionsverweis {offen}{}{}
und $\nabla$ ein
\definitionsverweis {linearer Zusammenhang}{}{}
auf dem trivialen
\definitionsverweis {Vektorbündel}{}{}

\mavergleichskette
{\vergleichskette
{ E
}
{ = }{ U \times \R^r
}
{  }{
}
{    }{
}
{   }{
}
}
{}{}{.}
Es seien

\mathbed {\partial_i} {}
 {1 \leq i \leq d} {}
 {} {} {} {}
die
\definitionsverweis {Standardvektorfelder}{}{}
und

\mathbed {s_j} {}
 {1 \leq j \leq r} {}
 {} {} {} {}
die Standardschnitte in $E$.}
\faktfolgerung {Dann gilt für den
\definitionsverweis {Krümmungsoperator}{}{}

\mavergleichskettedisp
{\vergleichskette
{ R( \partial_i , \partial_j)(s_k)
}
{ =} { \sum_{\ell = 1}^r  R^\ell_{ijk} s_\ell
}
{ } {
}
{ } {
}
{ } {
}
}
{}{}{}
mit

\mavergleichskettedisp
{\vergleichskette
{ R^\ell_{ijk}
}
{ =} { {\partial_i } \Gamma^\ell_{jk} - {\partial_j } \Gamma^\ell_{ik}  +  \sum_{ a  = 1}^r  \Gamma^a_{jk} \Gamma_{i a}^\ell -\sum_{ a  = 1}^r   \Gamma^a_{ik}  \Gamma_{j a}^\ell
}
{ } {
}
{ } {
}
{ } {
}
}
{}{}{,}
wobei die $\Gamma^\ell_{ik}$ die lokalen
\definitionsverweis {Christoffelsymbole}{}{}
des Zusammenhangs bezeichnen.}
\faktzusatz {}
\faktzusatz {}

}
{

Es ist

\mavergleichskettealignhandlinks
{\vergleichskettealignhandlinks
{ R( \partial_i ,\partial_j)
}
{ =} { \nabla_{\partial_i } \circ \nabla_{\partial_j } -  \nabla_{\partial_j } \circ \nabla_{\partial_i } - \nabla_{[\partial_i, \partial_j]}
}
{ =} { \nabla_{\partial_i } \circ \nabla_{\partial_j } -  \nabla_{\partial_j } \circ \nabla_{\partial_i }
}
{ } {
}
{ } {
}
}
{}
{}{,}
da ja die partiellen Ableitungen kommutieren und daher ihre Lie-Klammer gleich $0$ ist. Für einen Standardschnitt $s_k$ ist daher unter Verwendung von
[[Differenzierbare Mannigfaltigkeit/R/Vektorbündel/Linearer Zusammenhang/Lokale Beschreibung/Fakt|Lemma 25.10]]

\mavergleichskettealignhandlinks
{\vergleichskettealignhandlinks
{ R( \partial_i ,\partial_j) (s_k)
}
{ =} { \nabla_{\partial_i } ( \nabla_{\partial_j } (s_k) ) -  \nabla_{\partial_j } ( \nabla_{\partial_i } (s_k))
}
{ =} { \nabla_{\partial_i }  \left(  \sum_{\ell = 1}^r \Gamma^\ell_{jk} s_\ell  \right)  -  \nabla_{\partial_j }  \left(  \sum_{\ell = 1}^r \Gamma^\ell_{ik} s_\ell  \right)
}
{ =} { \sum_{\ell = 1}^r  \left(  \nabla_{\partial_i }  \left(  \Gamma^\ell_{jk} s_\ell  \right) - \nabla_{\partial_j }  \left(  \Gamma^\ell_{ik} s_\ell  \right)   \right)
}
{ =} { \sum_{\ell = 1}^r  \left(  {\partial_i } \Gamma^\ell_{jk} s_\ell +  \Gamma^\ell_{jk} \sum_{ \rho  = 1}^r \Gamma_{i \ell}^\rho  s_\rho  - {\partial_j } \Gamma^\ell_{ik} s_\ell - \Gamma^\ell_{ik} \sum_{ \rho  = 1}^r \Gamma_{j \ell}^\rho  s_\rho   \right)
}
}
{
\vergleichskettefortsetzungalign
{ =} { \sum_{\ell = 1}^r  \left(  {\partial_i } \Gamma^\ell_{jk} - {\partial_j } \Gamma^\ell_{ik}  \right) s_\ell  + \sum_{\ell = 1}^r \sum_{ \rho  = 1}^r  \left(  \Gamma^\ell_{jk} \Gamma_{i \ell}^\rho  s_\rho  - \Gamma^\ell_{ik}  \Gamma_{j \ell}^\rho  s_\rho  \right)
}
{ =} { \sum_{\ell = 1}^r  \left(  {\partial_i } \Gamma^\ell_{jk} - {\partial_j } \Gamma^\ell_{ik}  \right) s_\ell  + \sum_{\ell = 1}^r  \left(  \sum_{ a  = 1}^r  \Gamma^a_{jk} \Gamma_{i a}^\ell -\sum_{ a  = 1}^r   \Gamma^a_{ik}  \Gamma_{j a}^\ell  \right)  s_\ell
}
{ } {
}
{ } {}
}
{}{.}

}

__NOEDITSECTION__

\inputfaktbeweis
{Reelles Vektorbündel/Rang 1/Offene Menge/Krümmungsoperator/Fakt}
{Korollar}
{}
{

\faktsituation {Es sei

\mavergleichskette
{\vergleichskette
{ U
}
{ \subseteq }{ \R^n
}
{  }{
}
{    }{
}
{   }{
}
}
{}{}{}
eine
\definitionsverweis {offene Menge}{}{}
und sei $\nabla$ ein
\definitionsverweis {linearer Zusammenhang}{}{}
auf dem trivialen Vektorbündel
\mathl{U \times \R}{} vom
\definitionsverweis {Rang}{}{}
$1$ mit den
\definitionsverweis {Christoffelsymbolen}{}{}

\mathbed {\Gamma_i} {}
 {i = 1 , \ldots , n} {}
 {} {} {} {.}}
\faktuebergang {Dann gelten folgende Aussagen für den
\definitionsverweis {Krümmungsoperator}{}{}
zu den
\definitionsverweis {Standardvektorfeldern}{}{}
$\partial_i$.}
\faktfolgerung {\aufzaehlungzwei {Es ist

\mavergleichskettedisp
{\vergleichskette
{ R(\partial_i, \partial_j)
}
{ =} { \partial_i \Gamma_j - \partial_j \Gamma_i
}
{ } {
}
{ } {
}
{ } {
}
}
{}{}{.}
} {Genau dann ist

\mavergleichskette
{\vergleichskette
{ R(\partial_i, \partial_j)
}
{ = }{ 0
}
{  }{
}
{    }{
}
{   }{
}
}
{}{}{}
für alle $i,j$, wenn die
$1$-\definitionsverweis {Differentialform}{}{}
$\sum_{i = 1}^n \Gamma_i dx_i$
\definitionsverweis {geschlossen}{}{}
ist.
}}
\faktzusatz {}
\faktzusatz {}

}
{

\aufzaehlungzwei {Da der Rang des Bündels gleich $1$ ist, schreiben wir

\mavergleichskette
{\vergleichskette
{ \Gamma_i
}
{ = }{ \Gamma_{i1}^1
}
{  }{
}
{    }{
}
{   }{
}
}
{}{}{}
und entsprechend

\mavergleichskette
{\vergleichskette
{ R_{ij}
}
{ = }{ R_{ij1}^1
}
{  }{
}
{    }{
}
{   }{
}
}
{}{}{,}
die Aussage ist als

\mavergleichskettedisp
{\vergleichskette
{ R(\partial_i, \partial_j) (e)
}
{ =} { \left(  \partial_i \Gamma_j - \partial_j \Gamma_i  \right) e
}
{ } {
}
{ } {
}
{ } {
}
}
{}{}{}
mit dem konstanten Einheitsschnitt $e$ zu verstehen. Nach der expliziten Beschreibung in
[[Reelles Vektorbündel/Trivial/Linearer Zusammenhang/Krümmungsoperator/Beschreibung/Fakt|Lemma 28.2]]
ist

\mavergleichskettedisp
{\vergleichskette
{ R( \partial_i , \partial_j)
}
{ =} { {\partial_i } \Gamma_{j} - {\partial_j } \Gamma_{i}  + \Gamma_{j} \Gamma_{i } - \Gamma_{i}  \Gamma_{j }
}
{ =} { {\partial_i } \Gamma_{j} - {\partial_j } \Gamma_{i}
}
{ } {
}
{ } {
}
}
{}{}{.}
} {Dies folgt aus (1) und aus
[[Offene Teilmenge/R^n/Differentialform ersten Grades/Geschlossen/Aufgabe|Aufgabe 20.8]].
}

}

__NOEDITSECTION__

\inputfaktbeweis
{Reelles Vektorbündel/Linearer Zusammenhang/Krümmungsoperator/Eigenschaften/Fakt}
{Lemma}
{}
{

\faktsituation {Es sei $M$ eine zweifach stetige
\definitionsverweis {differenzierbare Mannigfaltigkeit}{}{}
und sei
\maabb {} { E } { M
} {}
ein zweifach stetig
\definitionsverweis {differenzierbares Vektorbündel}{}{}
über $M$ mit einem
\definitionsverweis {linearen Zusammenhang}{}{}
$\nabla$.}
\faktuebergang {Dann erfüllt der
\definitionsverweis {Krümmungsoperator}{}{}
die folgenden Eigenschaften.}
\faktfolgerung {\aufzaehlungdrei{Es ist

\mavergleichskettedisp
{\vergleichskette
{ R(V,W)
}
{ =} { -R(W,V)
}
{ } {
}
{ } {
}
{ } {
}
}
{}{}{.}
}{ $R$ ist additiv in beiden Komponenten.
}{Für eine differenzierbare Funktion
\maabb {f} { M } { \R
} {}
ist

\mavergleichskettedisp
{\vergleichskette
{ R(fV,W)
}
{ =} { f R(V,W)
}
{ } {
}
{ } {
}
{ } {
}
}
{}{}{.}
}}
\faktzusatz {}
\faktzusatz {}

}
{

Die Eigenschaften sind lokal in der Mannigfaltigkeit, man kann sie also auf einer offenen Überdeckung nachweisen, wobei die offenen Teilmengen diffeomorph zu offenen Mengen im $\R^d$ sind und worauf das Vektorbündel trivial ist. Somit folgt (1) aus
[[Reelles Vektorbündel/Trivial/Linearer Zusammenhang/Krümmungsoperator/Beschreibung/Fakt|Lemma 28.2]].
(2) folgt aus
[[Mannigfaltigkeit/Vektorbündel/R/Zusammenhang/Ableitung bezüglich Vektorfeld/Eigenschaften/Fakt|Lemma 24.10]]
und aus
[[Mannigfaltigkeit/Vektorfelder/Lie-Klammer/Eigenschaften/Fakt|Lemma 11.8]].
(3) Wir betrachten die lokale Situation mit den Vektorfeldern $f \partial_i$ und $g \partial_j$. Es ist

\mavergleichskettedisp
{\vergleichskette
{ R(f \partial_i, g \partial_j)
}
{ =} { \nabla_{f \partial_i } \circ \nabla_{g \partial_j } - \nabla_{g\partial_j } \circ \nabla_{f \partial_i } - \nabla_{[ f \partial_i, g \partial_j]}
}
{ } {
}
{ } {
}
{ } {
}
}
{}{}{.}
Nach
[[Mannigfaltigkeit/Vektorfelder/Lie-Klammer/Eigenschaften/Fakt|Lemma 11.8]]
ist

\mavergleichskettedisp
{\vergleichskette
{ [f \partial_i, g\partial_j]
}
{ =} { f [\partial_i, g\partial_j] - g \partial_j(f) \partial_i
}
{ } {
}
{ } {
}
{ } {
}
}
{}{}{.}
Somit ist für einen Schnitt $s$ unter Verwendung von
[[Mannigfaltigkeit/Vektorbündel/R/Linearer Zusammenhang/Leibnizregel/Fakt|Satz 25.5]]

\mavergleichskettealignhandlinks
{\vergleichskettealignhandlinks
{ R(f \partial_i, g\partial_j) (s)
}
{ =} { \nabla_{f \partial_i } \circ \nabla_{g \partial_j } (s)- \nabla_{g \partial_j } \circ \nabla_{f \partial_i } (s)- \nabla_{[ f \partial_i, g \partial_j]}(s)
}
{ =} { f  \nabla_{ \partial_i } ( \nabla_{g\partial_j } (s) ) - \nabla_{g\partial_j } ( \nabla_{f \partial_i } (s))  - \nabla_{ f [\partial_i, g\partial_j] - g \partial_j(f) \partial_i  } (s)
}
{ =} { f  \nabla_{ \partial_i } ( \nabla_{ g \partial_j } (s) )-\nabla_{ g \partial_j } (f \nabla_{ \partial_i } (s)) - \nabla_{f [\partial_i, g\partial_j] } (s)  + \nabla_{ g \partial_j (f) \partial_i }  (s)
}
{ =} { f  \nabla_{ \partial_i } ( \nabla_{ g \partial_j } (s) )- f \nabla_{ g \partial_j } ( \nabla_{ \partial_i } (s))- g \partial_j(f) \nabla_{\partial_i } (s)  - \nabla_{f [\partial_i, g\partial_j] } (s)          + g \partial_j( f)  \nabla_{\partial_i } (s)
}
}
{
\vergleichskettefortsetzungalign
{ =} { f \nabla_{ \partial_i } ( \nabla_{g \partial_j } (s) )- f \nabla_{ g \partial_j } ( \nabla_{ \partial_i } (s)) - f \nabla_{[\partial_i, g \partial_j]} (s)
}
{ =} { f R( \partial_i, \partial_j) (s)
}
{ } {
}
{ } {
}
}
{}{.}
Wegen (1) und (2) folgt daraus die Aussage.

}

__NOEDITSECTION__

\inputfaktbeweis
{Eindimensionale Mannigfaltigkeit/Reelles Vektorbündel/Linearer Zusammenhang/Krümmungsoperator trivial/Fakt}
{Korollar}
{}
{

\faktsituation {Es sei $M$ eine eindimensionale
$C^2$-\definitionsverweis {differenzierbare Mannigfaltigkeit}{}{}
und sei
\maabb {} { E } { M
} {}
ein zweifach stetig
\definitionsverweis {differenzierbares Vektorbündel}{}{}
über $M$ mit einem
\definitionsverweis {linearen Zusammenhang}{}{}
$\nabla$.}
\faktfolgerung {Dann ist der
\definitionsverweis {Krümmungsoperator}{}{}
für beliebige
\definitionsverweis {Vektorfelder}{}{}
$V,W$ trivial.}
\faktzusatz {}
\faktzusatz {}

}
{

Die Aussage ist lokal, daher können wir direkt

\mavergleichskette
{\vergleichskette
{ M
}
{ = }{ I
}
{  }{
}
{    }{
}
{   }{
}
}
{}{}{}
ein offenes Intervall und

\mavergleichskette
{\vergleichskette
{ V
}
{ = }{ f \partial
}
{  }{
}
{    }{
}
{   }{
}
}
{}{}{,}

\mavergleichskette
{\vergleichskette
{ W
}
{ = }{ g \partial
}
{  }{
}
{    }{
}
{   }{
}
}
{}{}{}
mit dem
\definitionsverweis {Standardvektorfeld}{}{}
$\partial$ ansetzen. Nach
[[Reelles Vektorbündel/Linearer Zusammenhang/Krümmungsoperator/Eigenschaften/Fakt|Lemma 28.4  (3)]]
ist

\mavergleichskettedisp
{\vergleichskette
{ R(f \partial, g \partial)
}
{ =} { fg R( \partial, \partial)
}
{ =} { 0
}
{ } {
}
{ } {
}
}
{}{}{}
wegen

\mavergleichskette
{\vergleichskette
{ [\partial, \partial ]
}
{ = }{ 0
}
{  }{
}
{    }{
}
{   }{
}
}
{}{}{.}

}

__NOEDITSECTION__

\inputfaktbeweis
{Reelles Vektorbündel/Trivialer Zusammenhang/Trivialer Krümmungsoperator/Fakt}
{Korollar}
{}
{

\faktsituation {Es sei $M$ eine
$C^2$-\definitionsverweis {differenzierbare Mannigfaltigkeit}{}{}
und

\mavergleichskette
{\vergleichskette
{ E
}
{ = }{ M \times \R^r
}
{  }{
}
{    }{
}
{   }{
}
}
{}{}{}
das
\definitionsverweis {triviale Vektorbündel}{}{}
mit dem
\definitionsverweis {trivialen Zusammenhang}{}{.}}
\faktfolgerung {Dann ist für beliebige
\definitionsverweis {Vektorfelder}{}{}
$V,W$ der
\definitionsverweis {Krümmungsoperator}{}{}
$R(V,W)$ trivial.}
\faktzusatz {}
\faktzusatz {}

}
{

Die Aussage ist lokal, wir können also von

\mavergleichskette
{\vergleichskette
{ M
}
{ = }{ U
}
{ \subseteq }{ \R^n
}
{    }{
}
{   }{
}
}
{}{}{}
offen ausgehen. Nach
[[Reelles Vektorbündel/Linearer Zusammenhang/Krümmungsoperator/Eigenschaften/Fakt|Lemma 28.4  (3)]]
können wir weiter auf den Fall reduzieren, dass die Vektorfelder
\definitionsverweis {Standardvektorfelder}{}{}
$\partial_i$ sind. Nach
[[R^n/Offene Menge/Trivialer Zusammenhang/Ableitungseigenschaften/Fakt/Beweis/Aufgabe|Aufgabe 25.11]]
sind die Christoffelsymbole zum trivialen Zusammenhang trivial, also folgt die Aussage aus
[[Reelles Vektorbündel/Trivial/Linearer Zusammenhang/Krümmungsoperator/Beschreibung/Fakt|Lemma 28.2]].

}

  \zwischenueberschrift{Der Satz von Frobenius}

Den folgenden Satz können wir nicht vollständig beweisen. Man beachte, dass man es bei
\zusatzklammer {lokal} {} {}
gegebenen Basisschnitten den Christoffelsymbolen nicht ansieht, ob bezüglich anderen Basisschnitten die Christoffelsymbole verschwinden. Dies wird eben durch das verschwinden der Krümmung charakterisiert.
__NOEDITSECTION__

\inputfaktbeweis
{Differenzierbare Mannigfaltigkeit/Vektorbündel/Linearer Zusammenhang/Lokal integrabel/Krümmung/Fakt}
{Satz}
{}
{

\faktsituation {Es sei $M$ eine
\definitionsverweis {differenzierbare Mannigfaltigkeit}{}{}
und
\maabbdisp {p} { E } { M
} {}
ein
\definitionsverweis {differenzierbares Vektorbündel}{}{,}
das mit einem
\definitionsverweis {linearen Zusammenhang}{}{}
$\nabla$ versehen sei.}
\faktuebergang {Dann sind folgende Aussagen äquivalent.}
\faktfolgerung {\aufzaehlungvier{Der Zusammenhang ist
\definitionsverweis {lokal integrabel}{}{.}
}{Der Zusammenhang ist lokal bezüglich geeigneter
\definitionsverweis {Basisschnitte}{}{}
isomorph zum
\definitionsverweis {trivialen Zusammenhang}{}{.}
}{Für jeden Punkt

\mavergleichskette
{\vergleichskette
{ P
}
{ \in }{ M
}
{  }{
}
{    }{
}
{   }{
}
}
{}{}{}
gibt es eine offene Kartenumgebung

\mavergleichskette
{\vergleichskette
{ P
}
{ \in }{ U
}
{  }{
}
{    }{
}
{   }{
}
}
{}{}{}
derart, dass $E \vert_U$ trivial ist und dass die zugehörigen
\definitionsverweis {Christoffelsymbole}{}{}
die Nullfunktionen sind.
}{Der
\definitionsverweis {Krümmungsoperator}{}{}
ist für beliebige
\definitionsverweis {Vektorfelder}{}{}
trivial.
}}
\faktzusatz {}
\faktzusatz {}

}
{

Alle Aussagen sind lokal, man kann also direkt von einer offenen Menge

\mavergleichskette
{\vergleichskette
{ U
}
{ \subseteq }{ \R^n
}
{  }{
}
{    }{
}
{   }{
}
}
{}{}{}
und dem trivialen Vektorbündel $U \times \R^r$ ausgehen. Sei (1) erfüllt. Dann kann man weiter annehmen, dass auf $U$ eine Familie von $r$ linear unabhängigen
\definitionsverweis {horizontalen Schnitten}{}{}

\mathl{s_1 , \ldots , s_r}{} gegeben ist. Nach
[[Mannigfaltigkeit/Vektorbündel/R/Zusammenhang/Global integrabel/Globale Trivialisierung/Aufgabe|Aufgabe 25.6]]
und
[[Mannigfaltigkeit/Vektorbündel/R/Linearer Zusammenhang/Global integrabel/Globale Trivialisierung des Zusammenhangs/Aufgabe|Aufgabe 25.7]]
kann man dann den Zusammenhang lokal trivialisieren. Dies ergibt (2). Die Äquivalenz von (2) und (3) beruht auf
[[R^n/Offene Menge/Trivialer Zusammenhang/Ableitungseigenschaften/Fakt/Beweis/Aufgabe|Aufgabe 25.11]]
und auf
[[Offene Menge/Triviales Vektorbündel/Linearer Zusammenhang/Spaltung/Christoffelsymbole/Bemerkung|Bemerkung 25.8]].
Wenn die Christoffelsymbole trivial sind, so sind die Standardschnitte horizontal und ergeben lokal horizontale Basisschnitte. Insgesamt sind also die ersten drei Formulierungen äquivalent. Wenn (2) gilt, so kann man lokal
[[Reelles Vektorbündel/Trivialer Zusammenhang/Trivialer Krümmungsoperator/Fakt|Korollar 28.6]]
anwenden und erhält, dass der Krümmungsoperator trivial ist.

Die Richtung von (4) nach (1) ist deutlich schwieriger.

}

Für Spezialfälle der Richtung von (4) nach (1) siehe
[[Eindimensionale Mannigfaltigkeit/Vektorbündel/R/Lokal integrabel/Aufgabe|Aufgabe 28.4]]
und
[[Differenzierbare Mannigfaltigkeit/Vektorbündel/R/Rang 1/Satz von Frobenius/Aufgabe|Aufgabe 28.5]].

 [[Kategorie:Latexseite]]
